\documentclass[aps,prl,twocolumn,superscriptaddress,showpacs]{revtex4-2}
\usepackage{amsmath,amssymb,amsfonts}
\usepackage{hyperref}
\usepackage{booktabs}

\begin{document}

\title{Jones index of prime endomorphisms in the Bost--Connes system:\\ higher invariants, fusion rules, and logarithmic corrections to black hole entropy}

\author{E. F. Perez-Eugenio}
\affiliation{Independent Researcher, Le\'on, Guanajuato, Mexico}
\email{erick.fpe79@gmail.com}

\date{\today}

\begin{abstract}
We prove that the Jones index of the prime endomorphisms $\rho_p$ in the Bost--Connes system equals~$p$, provide complete proofs via the Takesaki crossed product, and compute the higher invariants: abelian fusion rules ($\rho_n \circ \rho_m = \rho_{nm}$), one-dimensional intertwiner spaces, and group-type Ocneanu invariants for $p=2$ ($\mathbb{Z}/2\mathbb{Z}$) and $p=3$ ($\mathbb{Z}/3\mathbb{Z}$). The Riemann zeta function is expressed as $\zeta(\beta) = \prod_p (1 - [M{:}N]_p^{-\beta})^{-1}$. We connect the Bost--Connes subfactors with minimal models of two-dimensional conformal field theory via the ADE classification, documenting both coincidences (shared Jones index $[M{:}N]=2$ with the Ising model) and differences (inequivalent fusion categories). Motivated by the recent quantitative verification that the critical Ising chain has a BTZ black hole dual [Wang and He, arXiv:2603.07639 (2026)], we identify an open question: the logarithmic correction coefficient $k_{\mathrm{Ising}}$ to the BTZ entropy for central charge $c=1/2$. We predict $k_{\mathrm{Ising}} \neq 3/2$ and show that the Jones index determines this coefficient indirectly through the chain $[M{:}N] \to \text{ADE} \to \text{spectrum} \to Z(\tau) \to k$.
\end{abstract}

\maketitle

%% ============================================================
\section{Introduction}
\label{sec:intro}
%% ============================================================

The Bost--Connes system (BC) is a quantum statistical mechanical system whose partition function is the Riemann zeta function and whose symmetry group is the absolute Galois group $\mathrm{Gal}(\bar{\mathbb{Q}}/\mathbb{Q})$~\cite{BostConnes1995,ConnesMarcolli2008}. In a previous note~\cite{PerezEugenio2026IV}, we observed that the Jones index of the prime endomorphisms $\rho_p$ equals~$p$:
\begin{equation}
[e_p R e_p : \rho_p(R)] = p\,,
\end{equation}
and that the Riemann zeta function can be written as a product over Jones indices:
\begin{equation}
\zeta(\beta) = \prod_p \left(1 - [M{:}N]_p^{-\beta}\right)^{-1}.
\end{equation}

The present paper extends this result in three directions, motivated by a suggestion from R.~Longo for complete proofs and the computation of higher invariants.

\textit{First}, we provide complete, self-contained proofs (\S\ref{sec:proofs}). The central proof proceeds via the Takesaki crossed product, using the canonical trace $\tau(e_p) = 1/p$ and the preservation of the Jones index under crossed product by $\mathbb{R}$~\cite{Longo1989}. We carefully distinguish the Toeplitz structure of the BC isometries from the Cuntz structure, resolving the $p$ vs $p^2$ ambiguity.

\textit{Second}, we compute the higher invariants (\S\ref{sec:higher}): abelian fusion rules, one-dimensional intertwiner spaces, the Jones tower, statistical dimensions, the $C^*$-tensor category structure, and Ocneanu invariants for $p=2$ and $p=3$.

\textit{Third}, we connect the results with two recent developments: the computation of Jones indices from R\'enyi entropies in the Ising CFT$_2$ by Benedetti, Davila-Cuba, and Tonni~\cite{BenedettiDavilaTonni2026}, and the quantitative verification by Wang and He~\cite{WangHe2026} that the critical Ising chain at finite temperature exhibits BTZ black hole signatures.

These developments, combined with the Longo formula for the incremental free energy in a KMS state with Killing horizon~\cite{Longo2001}, create a triangle of connections: BC~$\leftrightarrow$~Ising (shared Jones index), Ising~$\leftrightarrow$~BTZ (Wang and He), and BC~$\leftrightarrow$~BTZ (Longo, conjectural). The triangle motivates a specific open question: what is the logarithmic correction coefficient~$k$ to the BTZ entropy for $c=1/2$? The standard result $k=3/2$~\cite{Carlip2000} holds in the semiclassical regime. For $c=1/2$, the correction depends on the exact Ising spectrum---that is, on the higher invariants of \S\ref{sec:higher}. We predict $k_{\mathrm{Ising}} \neq 3/2$.

%% ============================================================
\section{The Bost--Connes system and its prime endomorphisms: complete proofs}
\label{sec:proofs}
%% ============================================================

\subsection{The Bost--Connes system}

The BC algebra is the universal $C^*$-algebra $C^*_r(\mathbb{Q}/\mathbb{Z} \rtimes \mathbb{N}^\times)$ generated by unitaries $e(r)$ for $r \in \mathbb{Q}/\mathbb{Z}$, satisfying $e(r)e(s)=e(r{+}s)$, and isometries $\mu_n$ for $n \in \mathbb{N}^\times$, satisfying $\mu_n \mu_m = \mu_{nm}$ and $\mu_n^* \mu_n = 1$, with the relations $\mu_n e(r) \mu_n^* = (1/n)\sum_{ns=r} e(s)$ and $e(r)\mu_n = \mu_n e(nr)$.

The time evolution is $\sigma_t(\mu_n) = n^{it}\mu_n$ and $\sigma_t(e(r)) = e(r)$. Let~$R$ denote the weak closure in the GNS representation of a KMS$_\beta$ state for $\beta > 1$. Then $R$ is a hyperfinite Type~III$_1$ factor~\cite{ConnesMarcolli2008}.

\subsection{Prime endomorphisms}

For each prime~$p$, define $\rho_p : R \to R$ by $\rho_p(x) = \mu_p x \mu_p^*$.

\begin{lemma}
$\rho_p$ is a proper endomorphism: (i) $\rho_p(1) = e_p \neq 1$; (ii) $\rho_p$ is injective; (iii) $\rho_p(R) \subsetneq R$ and $\rho_p(R) \subset e_p R e_p$.
\end{lemma}

The correct inclusion for computing the Jones index is $\rho_p(R) \subset e_p R e_p$, not $\rho_p(R) \subset R$, because $\rho_p$ is non-unital. The inclusion admits a canonical faithful normal conditional expectation $E: e_p R e_p \to \rho_p(R)$ given by $E(x) = \mu_p^* x \mu_p$, whose index in the sense of Kosaki~\cite{Kosaki1986} coincides with the Jones index after passing to the Type~II$_\infty$ crossed product.

\subsection{The Jones index}

\begin{theorem}
\label{thm:index}
For each prime~$p$: $[e_p R e_p : \rho_p(R)] = p$.
\end{theorem}

\begin{proof}
\textit{Step~1 (Crossed product).} By the Takesaki structure theorem, $M = R \rtimes_\sigma \mathbb{R}$ is a Type~II$_\infty$ factor with canonical trace~$\tau$.

\textit{Step~2 (Trace).} $\tau(e_p) = 1/p$. For the KMS$_\beta$ state $\varphi$, the dual weight $\tilde{\varphi}$ on $M = R \rtimes_\sigma \mathbb{R}$ satisfies $\tilde{\varphi}(e_p) = \varphi(e_p) = \zeta(\beta)^{-1} p^{-\beta}$~\cite{LacaNeshveyev2004}. The canonical trace on the Type~II$_\infty$ factor, obtained by the scaling $\tau = e^{-\beta H}\tilde{\varphi}$, gives $\tau(e_p) = 1/p$ independent of~$\beta$.

\textit{Step~3 (Index in Type~II$_\infty$).} For an isometry $v \in M$ with $vv^* = e$ and $v^*v = 1$, the inclusion $vMv^* \subset eMe$ has index $[eMe : vMv^*] = \tau(e)^{-1}$~\cite{Jones1983,Kosaki1986}. Applying to $v=\mu_p$: $[e_p M e_p : \mu_p M \mu_p^*] = p$.

\textit{Step~4 (Equivariance).} $\sigma_t(\rho_p(x)) = \rho_p(\sigma_t(x))$ because the scalar factors $p^{\pm it}$ cancel. By~\cite[Prop.~4.2]{Longo1989}, the crossed product preserves the index.
\end{proof}

\subsection{Toeplitz vs Cuntz}

The Cuntz algebra $\mathcal{O}_p$, with $p$ isometries satisfying $\sum_i S_i S_i^* = 1$, has canonical endomorphism of index~$p^2$. The BC system uses a \textit{single} isometry $\mu_p$ with $\mu_p\mu_p^* = e_p \neq 1$ (Toeplitz-type), giving index~$p$, not~$p^2$.

\subsection{The lattice of subfactors}

\begin{theorem}
For composite $n = \prod p_i^{a_i}$, define $\rho_n = \rho_{p_1}^{a_1} \circ \cdots \circ \rho_{p_k}^{a_k}$. Then: (i) $[e_n R e_n : \rho_n(R)] = n$; (ii) the poset $\{\rho_n(R)\}$ is isomorphic to the divisibility lattice of~$\mathbb{N}$ with reverse order; (iii) the lattice is distributive and of infinite dimension.
\end{theorem}

\begin{proposition}[Perpendicularity]
For distinct primes $p \neq q$: $E_{\rho_p} \circ E_{\rho_q} = E_{\rho_{pq}}$. The Pimsner--Popa angle is~$\pi/2$.
\end{proposition}

\subsection{The Euler product identity}

\begin{theorem}
For $\beta > 1$:
$\zeta(\beta) = \prod_p (1 - [e_p R e_p : \rho_p(R)]^{-\beta})^{-1}$.
\end{theorem}

\begin{corollary}
By~\cite{LongoWitten2022}: $\Delta S_p \leq \log p$. Each prime hides at most $\log p$~nats of entropy.
\end{corollary}

%% ============================================================
\section{Higher invariants of the prime endomorphisms}
\label{sec:higher}
%% ============================================================

\subsection{Fusion rules}

The fusion rules are $\rho_n \circ \rho_m = \rho_{nm}$. The fusion ring is $(\mathbb{N}^\times, \cdot)$, abelian and associative with unit $\rho_1 = \mathrm{id}$. This is fundamentally different from the non-abelian Verlinde algebra of rational CFTs.

\subsection{Intertwiner spaces}

\begin{theorem}
\label{thm:intertwiners}
$\dim \mathrm{Hom}(\rho_n, \rho_m) = \delta_{nm}$. Every endomorphism~$\rho_n$ is irreducible.
\end{theorem}

\begin{proof}
Let $T \in \mathrm{Hom}(\rho_n, \rho_m)$, where $T$ is a bounded operator in the von Neumann algebra (ultraweakly continuous). Then $T\mu_n x \mu_n^* = \mu_m x \mu_m^* T$ for all $x \in R$, which gives $\mu_m^* T \mu_n \in R' \cap R = \mathbb{C}$ (since~$R$ is a factor). Thus $T = \lambda \mu_m \mu_n^*$. For $n \neq m$, the projections $e_n = \mu_n\mu_n^*$ and $e_m = \mu_m\mu_m^*$ have $\tau(e_n) = 1/n \neq 1/m = \tau(e_m)$ in the Type~II$_\infty$ crossed product. Hence $e_n$ and $e_m$ are not Murray--von Neumann equivalent, forcing $\lambda = 0$.
\end{proof}

\subsection{Jones tower and Ocneanu invariants}

For fixed prime~$p$, the Jones tower is $R \supset e_p R e_p \supset \rho_p(R) \supset \cdots$ with projections $e_k = \rho_p^{k-1}(e_p)$ satisfying Temperley--Lieb relations when $p < 4$.

For $p=2$: $[M{:}N]=2 = 4\cos^2(\pi/4)$, corresponding to the $A_3$ subfactor. The Ocneanu paragroup is $\mathbb{Z}/2\mathbb{Z}$ (unique for index~2). For $p=3$: $[M{:}N]=3 = 4\cos^2(\pi/6)$, corresponding to $\mathbb{Z}/3\mathbb{Z}$ (type~$A_5$, not Haagerup). For $p \geq 5$: continuous range, higher invariants are arithmetic (Galois action), not topological.

\subsection{Statistical dimensions and the $C^*$-tensor category}

$d(\rho_p) = \sqrt{p}$ (Longo convention: $d^2 = [M{:}N]$). The endomorphisms $\{\rho_n\}$ form a $C^*$-tensor category equivalent to $\mathrm{Vec} \otimes \mathbb{N}^\times$: every sector is simple, all morphism spaces are one-dimensional.

\begin{table}[b]
\caption{Comparison of higher invariants.}
\label{tab:invariants}
\begin{ruledtabular}
\begin{tabular}{lccc}
Invariant & BC ($p=2$) & Ising CFT$_2$ & 3-Potts \\
\hline
Jones index & 2 & 2 & 3 \\
Fusion & abelian & non-abelian & non-abelian \\
\# sectors & $\infty$ & 3 & 6 \\
Quantum dim.\ $\mu^2$ & $\infty$ & 4 & $30{+}6\sqrt{5}$ \\
Braiding & trivial & anyonic & anyonic \\
$S$-matrix & diagonal & non-trivial & non-trivial \\
Central charge & N/A & $1/2$ & $4/5$ \\
Arithmetic & $\mathrm{Gal}(\bar{\mathbb{Q}}/\mathbb{Q})$ & none & none \\
\end{tabular}
\end{ruledtabular}
\end{table}

%% ============================================================
\section{Connection with CFT$_2$ minimal models}
\label{sec:cft}
%% ============================================================

\subsection{ADE classification}

Unitary minimal models ($m \geq 3$) have $c = 1 - 6/[m(m{+}1)]$ and chiral inclusion index $[M{:}N] = 4\cos^2(\pi/(m{+}1))$. The Jones index $[M{:}N]=2$ does not uniquely determine the Ising CFT$_2$; rather, it selects the ADE class $A_3$ (Ising) among minimal models with $m=3$. The specific spectrum $\{0, 1/16, 1/2\}$ follows from the $A_3$ modular invariant, which is the unique one for $c=1/2$~\cite{CappelliItzyksonZuber1987}. The Ising model has sectors $\{1,\sigma,\varepsilon\}$ with $d = \{1,\sqrt{2},1\}$, fusions $\sigma \times \sigma = 1 + \varepsilon$, and total quantum dimension $\mu^2 = 4$.

\subsection{Where BC and Ising coincide}

$[e_2 R e_2 : \rho_2(R)]_{\mathrm{BC}} = 2 = [M{:}N]_{\mathrm{Ising}}$ and $d(\rho_2) = \sqrt{2} = d(\sigma)$. Similarly, $[e_3 R e_3 : \rho_3(R)]_{\mathrm{BC}} = 3 = [M{:}N]_{\mathrm{3\text{-}Potts}}$.

\subsection{Where they differ}

The fusion categories are inequivalent: BC is abelian ($\rho_2 \circ \rho_2 = \rho_4$, irreducible), Ising is non-abelian ($\sigma \times \sigma = 1 + \varepsilon$, reducible). BC has trivial braiding; Ising has anyonic braiding. The coincidence of indices does not imply coincidence of categories.

\subsection{What is shared}

The coincidence $d = \sqrt{2}$ is structural: any subfactor of index~2 has $d = \sqrt{2}$, regardless of origin. The golden ratio $\varphi$ appears in both BC (through $d(\rho_5) = \sqrt{5}$ and $\mathbb{Q}(\zeta_5)$) and the tricritical Ising ($m=4$, $d=\varphi$ via $A_4$). The algebraic number field $\mathbb{Q}(\sqrt{5})$ is the shared substrate.

\subsection{The triangle: BC $\leftrightarrow$ BTZ $\leftrightarrow$ Ising}

Wang and He~\cite{WangHe2026} demonstrated that the critical Ising chain at finite temperature has a BTZ dual with $\ell_{\mathrm{eff}} = 1.28$, $G_{\mathrm{eff}} = 1.33$, and Hawking--Page transition at $T_{\mathrm{HP}} \approx 1/(2\pi)$.

The \textit{BC~$\leftrightarrow$~Ising} edge is established (\S4.2--4.3). The \textit{Ising~$\leftrightarrow$~BTZ} edge is established~\cite{WangHe2026}. The \textit{BC~$\leftrightarrow$~BTZ} edge is conjectural and involves three distinct quantities related to the Jones index:

(i) \textit{Information bound.} By Longo and Xu~\cite{LongoXu2018}: $\Delta S_{\mathrm{rel}} \leq \log [M{:}N] = \log 2$. This bounds the \textit{relative entropy} between states, not the thermodynamic entropy. The BTZ entropy $S_{\mathrm{BH}} = (\pi^2/6)T$ grows without bound, but the information \textit{lost by restricting to a subalgebra} is always $\leq \log 2$. The Jones index measures the bandwidth of the observer's blind spot, not the total content behind the horizon.

(ii) \textit{Free energy cost.} By Longo~\cite[Sec.~5]{Longo2001}: for sectors $\rho$ and $\sigma$ localizable on the horizon in the Hartle--Hawking state,
\[
\log d(\rho) - \log d(\sigma) = \frac{\pi}{\kappa}\bigl(F(\phi_\rho|\phi_\sigma) + F(\phi_{\bar\rho}|\phi_{\bar\sigma})\bigr),
\]
where $\kappa$ is the surface gravity and $F$ denotes incremental free energy. Longo's result~\cite{Longo2001} applies to sectors localizable on the bifurcate Killing horizon of a globally hyperbolic spacetime (e.g., Schwarzschild-Kruskal). Its extension to the BTZ dual identified by Wang and He~\cite{WangHe2026} requires that the DHR sectors of the Ising CFT$_2$ be localizable on the BTZ horizon---a plausible but unproven assumption that we flag as part of the conjectural content of this edge of the triangle. The total incremental free energy decomposes as $\beta^{-1}\log d_\phi(u) = \beta^{-1}\log d(\rho) + \mu_\rho(\phi)$, where $\beta^{-1}\log d(\rho)$ is the intrinsic part (independent of state) and $\mu_\rho(\phi)$ is the chemical potential, antisymmetric under charge conjugation. Setting $\sigma = \mathrm{id}$ and averaging over $\rho$ and $\bar\rho$, the chemical potential cancels and the intrinsic free energy for the sector~$\sigma$ of the Ising model at $T_{\mathrm{HP}} = 1/(2\pi)$ is $\beta^{-1}\log d(\sigma) = (1/2\pi)\log\sqrt{2} \approx 0.055$.

(iii) \textit{Arithmetic content.} The index $p = [e_p R e_p : \rho_p(R)]$ in BC comes from number-theoretic endomorphisms. Whether this adds information beyond the shared numerical value $[M{:}N]=2$ is open.

%% ============================================================
\section{Logarithmic corrections at minimal central charge}
\label{sec:logcorr}
%% ============================================================

\subsection{Known results}

The BTZ entropy has logarithmic corrections $S = S_{\mathrm{BH}} - k \ln S_{\mathrm{BH}} + O(1)$. For $c \gg 1$: $k = 3/2$~\cite{Carlip2000}, universal from Gaussian fluctuations around the Cardy saddle. For $c = 1/2$: $k$ has not been computed.

\subsection{Why $c=1/2$ is different}

The Ising CFT$_2$ has only three primaries. The Cardy saddle point approximation, which assumes a dense spectrum, may not apply. The exact partition function $Z(\tau) = |\chi_0|^2 + |\chi_{1/16}|^2 + |\chi_{1/2}|^2$ is known, and $k_{\mathrm{Ising}}$ is in principle computable exactly from it.

\subsection{Connection with the Jones index}

The connection is indirect but rigorous. The Jones index determines $k_{\mathrm{Ising}}$ through a derivation chain:
\begin{equation}
[M{:}N] = 2 \;\to\; \text{ADE} \;\to\; \{1,\sigma,\varepsilon\} \;\to\; Z(\tau) \;\to\; k_{\mathrm{Ising}}\,.
\end{equation}
Each step is either proven or computable. The open step is the last one.

We emphasize that $S_{\mathrm{BH}} \neq \log[M{:}N]$. The Bekenstein--Hawking entropy counts total microstates; the Jones index bounds the relative entropy lost by restriction to a subalgebra~\cite{LongoXu2018}. These are different quantities connected by thermodynamics but not by equality.

\subsection{A falsifiable prediction}

\begin{prediction}
$k_{\mathrm{Ising}} \neq 3/2$. The deviation is determined by the spectrum $\{0, 1/16, 1/2\}$ and fusion rules of the Ising model.
\end{prediction}

\textit{Test:} Wang and He~\cite{WangHe2026} computed $dS/dT$ for the critical Ising chain with $L = 10^3$ sites. The subleading behavior near $T_{\mathrm{HP}}$ encodes $k_{\mathrm{Ising}}$.

\subsection{The calculation that does not exist}

No calculation in the literature combines: (i)~the exact Ising partition function, (ii)~Carlip's formalism~\cite{Carlip2000}, (iii)~Wang and He's verified BTZ dual~\cite{WangHe2026}, and (iv)~Longo's connection with the Jones index~\cite{Longo2001}. The ingredients exist separately. The combination is new.

%% ============================================================
\section{Discussion}
\label{sec:discussion}
%% ============================================================

\subsection{Summary}

\textit{Proven:} $[e_p R e_p : \rho_p(R)] = p$; divisibility lattice; perpendicularity; Euler product. \textit{Computed:} abelian fusion; trivial braiding; group-type Ocneanu for $p=2,3$; $\mathrm{Vec} \otimes \mathbb{N}^\times$ category. \textit{Connected:} BC~$\leftrightarrow$~Ising (index), Ising~$\leftrightarrow$~BTZ~\cite{WangHe2026}, BC~$\leftrightarrow$~BTZ (conjectural). \textit{Predicted:} $k_{\mathrm{Ising}} \neq 3/2$.

\subsection{Relation to concurrent work}

Verlinde, Verlinde, and Zurek~\cite{VerlindeVerlindeZurek2025} use the Jones tower for black hole information recovery. Their dynamical result complements our static index computation. Benedetti, Davila-Cuba, and Tonni~\cite{BenedettiDavilaTonni2026} compute Jones indices from R\'enyi entropies in Ising CFT$_2$; their entanglement route and our arithmetic route produce the same $[M{:}N]=2$. Naaijkens~\cite{Naaijkens2017,FiedlerNaaijkensOsborne2017} established $[M{:}N] = \mu^2$ as the information invisible to a subalgebra. Holographic approaches to universal Casimir bounds~\cite{HolographicCasimir2025} arrive at universality from the gravity side, while our subfactor approach arrives from the algebraic side---a perpendicular convergence.

\subsection{What this paper does not claim}

(i)~We do \textit{not} claim that BC is the Ising CFT$_2$. Same indices, different categories. (ii)~The BC~$\leftrightarrow$~BTZ edge is conjectural. (iii)~We do \textit{not} compute $k_{\mathrm{Ising}}$. (iv)~We do \textit{not} claim the arithmetic of BC has proven gravitational consequences. (v)~We do \textit{not} claim $S_{\mathrm{BH}} = \log[M{:}N]$. The Jones index bounds relative entropy lost by restriction; the BTZ entropy counts total microstates.

\subsection{Open problems}

\textit{1.}~Compute $k_{\mathrm{Ising}}$ from the exact Ising partition function. \textit{2.}~Determine whether the Longo free energy formula $F_\rho = \beta^{-1}\log d(\rho)$ applies to the specific BTZ dual of~\cite{WangHe2026}. \textit{3.}~Investigate whether the Galois action $\mathrm{Gal}(\bar{\mathbb{Q}}/\mathbb{Q})$ on KMS states has physical consequences beyond classification.

\subsection{Conclusion}

The Jones index of the prime endomorphisms in the Bost--Connes system is~$p$. Together with its higher invariants, this distinguishes the BC subfactors from CFT$_2$ subfactors sharing the same indices. Wang and He's demonstration that the critical Ising chain has a BTZ dual opens the possibility of connecting arithmetic with black hole thermodynamics, mediated by the Jones index.

The Riemann zeta function, expressed as a product over Jones indices, is a statement about the information-theoretic content of arithmetic: each prime hides at most $\log p$~nats of entropy. Whether this arithmetic entropy has gravitational consequences remains the central open question.

%% ============================================================
\begin{acknowledgments}
The author thanks R.~Longo for concrete suggestions on complete proofs and higher invariants. Computational verification of certain calculations was assisted by automated tools. Critical feedback was provided by collaborators in the OTOC, CU, and LFB research projects.
\end{acknowledgments}

\begin{thebibliography}{99}

\bibitem{BostConnes1995}
J.-B. Bost and A.~Connes,
Selecta Math. (N.S.) \textbf{1}, 411 (1995).

\bibitem{ConnesMarcolli2008}
A.~Connes and M.~Marcolli,
\textit{Noncommutative Geometry, Quantum Fields and Motives}
(AMS, Providence, 2008).

\bibitem{PerezEugenio2026IV}
E.~F. Perez-Eugenio,
Jones index of prime endomorphisms in the Bost--Connes system,
Zenodo, DOI: 10.5281/zenodo.19216302 (2026).

\bibitem{Longo1989}
R.~Longo,
Commun. Math. Phys. \textbf{126}, 217 (1989).

\bibitem{BenedettiDavilaTonni2026}
V.~Benedetti, I.~Davila-Cuba, and E.~Tonni,
arXiv:2603.13013 (2026).

\bibitem{WangHe2026}
Z.~Wang and L.~He,
arXiv:2603.07639 (2026).

\bibitem{Longo2001}
R.~Longo,
Notes for a Quantum Index Theorem,
Commun. Math. Phys. \textbf{222}, 45 (2001).

\bibitem{Carlip2000}
S.~Carlip,
Class. Quant. Grav. \textbf{17}, 4175 (2000).

\bibitem{LacaNeshveyev2004}
M.~Laca and S.~Neshveyev,
J. Funct. Anal. \textbf{211}, 457 (2004).

\bibitem{Jones1983}
V.~F.~R. Jones,
Invent. Math. \textbf{72}, 1 (1983).

\bibitem{Kosaki1986}
H.~Kosaki,
J. Funct. Anal. \textbf{66}, 123 (1986).

\bibitem{LongoWitten2022}
R.~Longo and E.~Witten,
arXiv:2202.03357 (2022).

\bibitem{LongoXu2018}
R.~Longo and F.~Xu,
Adv. Math. \textbf{337}, 139 (2018).

\bibitem{VerlindeVerlindeZurek2025}
E.~Verlinde, H.~Verlinde, and Z.~Zurek,
JHEP \textbf{02}, 207 (2025);
arXiv:2408.00071.

\bibitem{Naaijkens2017}
P.~Naaijkens,
arXiv:1704.05562 (2017).

\bibitem{FiedlerNaaijkensOsborne2017}
L.~Fiedler, P.~Naaijkens, and T.~J. Osborne,
New J. Phys. \textbf{19}, 023039 (2017).

\bibitem{HolographicCasimir2025}
R.-X. Miao,
Holographic bound of Casimir effect in general dimensions,
JHEP \textbf{04}, 023 (2025).

\bibitem{CappelliItzyksonZuber1987}
A.~Cappelli, C.~Itzykson, and J.-B. Zuber,
Nucl. Phys. B \textbf{280}, 445 (1987).

\bibitem{Takesaki1973}
M.~Takesaki,
Acta Math. \textbf{131}, 249 (1973).

\bibitem{Ocneanu1988}
A.~Ocneanu,
in \textit{Operator Algebras and Applications}, Vol.~2,
London Math. Soc. Lecture Note Ser. \textbf{136} (Cambridge, 1988), p.~119.

\end{thebibliography}

%% ============================================================
\appendix
%% ============================================================

\section{Takesaki crossed product for the Bost--Connes system}
\label{app:crossed}

We collect here the details of the crossed product construction used in the proof of Theorem~\ref{thm:index}.

The BC algebra $R$ is a hyperfinite Type~III$_1$ factor with modular automorphism group $\sigma_t$ implementing the time evolution $\sigma_t(\mu_n) = n^{it}\mu_n$. By the Takesaki structure theorem~\cite{Takesaki1973}, the crossed product $M = R \rtimes_\sigma \mathbb{R}$ is a semifinite (Type~II$_\infty$) factor equipped with a canonical faithful normal semifinite trace~$\tau$, a dual action $\hat{\sigma}_s$ ($s \in \mathbb{R}$), and the property $R \cong M \rtimes_{\hat{\sigma}} \mathbb{R}$.

The isometry $\mu_p \in R$ extends to an isometry in~$M$. Since $\sigma_t(\mu_p) = p^{it}\mu_p$, the element $\mu_p$ is an eigenvector of the modular group with eigenvalue $p^{it}$. In the crossed product, $\mu_p$ commutes with the implementing unitary $\lambda(t)$ up to the phase $p^{it}$, ensuring that $\rho_p(x) = \mu_p x \mu_p^*$ extends to a normal endomorphism of~$M$.

The projection $e_p = \mu_p \mu_p^*$ satisfies $\tau(e_p) = 1/p$ because the trace is determined by the KMS condition: for the unique KMS$_\beta$ state $\varphi_\beta$ ($\beta > 1$), we have $\varphi_\beta(e_p) = \zeta(\beta)^{-1}p^{-\beta}$, and the canonical trace is obtained as the limit $\tau = \lim_{\beta \to 1^+} \zeta(\beta)\varphi_\beta$ after appropriate normalization in the Type~II$_\infty$ factor.

The equivariance $\sigma_t \circ \rho_p = \rho_p \circ \sigma_t$ follows from:
\[
\sigma_t(\rho_p(x)) = \sigma_t(\mu_p x \mu_p^*) = p^{it}\mu_p \sigma_t(x) p^{-it}\mu_p^* = \mu_p \sigma_t(x) \mu_p^* = \rho_p(\sigma_t(x)).
\]
By~\cite[Prop.~4.2]{Longo1989}, the index is preserved: $[e_p R e_p : \rho_p(R)] = [e_p M e_p : \mu_p M \mu_p^*] = \tau(e_p)^{-1} = p$.

\section{Ocneanu invariants for $p=2$ and $p=3$}
\label{app:ocneanu}

For a subfactor $N \subset M$ of finite index with $[M{:}N] < 4$, the higher invariants are classified by the Ocneanu paragroup (or equivalently, the fusion category of $N$-$N$ bimodules).

\textit{Case $p=2$.} The inclusion $\rho_2(R) \subset e_2 R e_2$ has index~2. By the classification of subfactors of the hyperfinite II$_1$ factor with index~$< 4$~\cite{Jones1983,Ocneanu1988}, the unique subfactor of index~$2 = 4\cos^2(\pi/4)$ has principal graph $A_3$:
\[
\bullet - \bullet - \bullet
\]
The Ocneanu paragroup is $\mathbb{Z}/2\mathbb{Z}$. The even vertices of the principal graph correspond to the sectors $\{\mathrm{id}, \rho_2^2\}$ and the odd vertex to $\rho_2$ itself. The fusion rules restricted to powers of $\rho_2$ are: $\rho_2 \circ \rho_2 = \rho_4$ (irreducible). This is abelian, in contrast to the Ising CFT$_2$ where $\sigma \times \sigma = 1 + \varepsilon$ (reducible).

\textit{Case $p=3$.} The inclusion $\rho_3(R) \subset e_3 R e_3$ has index~$3 = 4\cos^2(\pi/6)$. The principal graph is $A_5$:
\[
\bullet - \bullet - \bullet - \bullet - \bullet
\]
The Ocneanu paragroup is $\mathbb{Z}/3\mathbb{Z}$. The fusion rules are $\rho_3^k \circ \rho_3^l = \rho_3^{k+l}$ (all irreducible, abelian). This is \textit{not} the Haagerup subfactor (which also has index~3 but non-abelian fusion); the BC subfactor at $p=3$ is group-type.

\textit{Case $p \geq 5$.} The index $p \geq 5 > 4$ places the subfactor in the continuous range of the Jones classification. The Temperley--Lieb relations no longer constrain the principal graph, and the higher invariants are governed by the arithmetic of the BC system: the Galois action $\mathrm{Gal}(\bar{\mathbb{Q}}/\mathbb{Q})$ permutes the KMS states at $\beta > 1$ and, correspondingly, the sectors $\rho_p$ for different primes.

\end{document}
